\Needspace{9\baselineskip}\begin{center}\textit{Umwandlung gewöhnlicher Brüche in Dezimalbrüche.}\end{center}

\Needspace{7\baselineskip}
\Article{312}
\textsc{Definition.} Wird ein gewöhnlicher Bruch in einen Dezimalbruch umgewandelt, so nennen wir die Folge der Dezimalziffern\HistoricalNote{*)}{Der Kürze wegen beschränken wir die folgende Untersuchung auf das gewöhnliche Dezimalsystem, da sie sich leicht auf jedes andere ausdehnen lässt.} ohne den etwa vorhandenen ganzzahligen Teil, ob sie endlich ist oder unendlich fortläuft, die \textit{Mantisse} des Bruchs. Wir gebrauchen diesen sonst nur bei Logarithmen üblichen Ausdruck in einem weiteren Sinn. So ist \textit{zum Beispiel} die Mantisse des Bruchs \(\frac18\) gleich \(125\), die von \(\frac{35}{16}\) gleich \(1875\) und die von \(\frac2{37}\) gleich \(054054\ldots\), ins Unendliche fortgesetzt.

Aus dieser Definition folgt sofort, dass die Brüche \(\frac ln,\frac mn\) mit demselben Nenner gleiche oder verschiedene Mantissen haben, je nachdem ihre Zähler \(l,m\) nach \(n\) kongruent oder inkongruent sind. Eine endliche Mantisse ändert sich nicht, wenn rechts beliebig viele Nullen angefügt werden. Die Mantisse des Bruchs \(\frac{10m}{n}\) erhält man, indem man die erste Ziffer der Mantisse von \(\frac mn\) entfernt. Allgemein findet man die Mantisse von \(\frac{10^\nu m}{n}\), indem man die ersten \(\nu\) Ziffern der Mantisse von \(\frac mn\) entfernt. Die Mantisse von \(\frac1n\) beginnt sofort mit einer signifikanten Ziffer, \textit{das heißt}, einer von null verschiedenen, wenn \(n\) nicht \(>10\) ist. Gilt dagegen \(n>10\), ist der Nenner keiner Potenz von \(10\) gleich und besteht er aus \(k\) Ziffern, so sind die ersten \(k-1\) Ziffern der Mantisse von \(\frac1n\) Nullen, und erst die folgende, die \(k\)-te, ist signifikant. Daraus folgt leicht: Haben \(\frac ln,\frac mn\) verschiedene Mantissen, \textit{das heißt},
\SourceLeaf{388}
sind \(l,m\) nach \(n\) inkongruent, so können diese Mantissen sicher nicht in den ersten \(k\) Ziffern übereinstimmen, sondern müssen sich spätestens in der \(k\)-ten Ziffer unterscheiden.

\Needspace{7\baselineskip}
\Article{313}
\textsc{Aufgabe.} \textit{Aus dem gegebenen Nenner eines Bruchs \(\frac mn\) und den ersten \(k\) Ziffern seiner Mantisse den Zähler \(m\) finden, den wir kleiner als \(n\) voraussetzen.}

\textit{Lösung.} Man betrachte diese \(k\) Ziffern als ganze Zahl, multipliziere sie mit \(n\) und dividiere das Produkt durch \(.10^k\), oder entferne die letzten \(k\) Ziffern. Ist der Quotient eine ganze Zahl, sind also die entfernten Ziffern Nullen, so ist er offenbar der gesuchte Zähler und die gegebene Mantisse vollständig. Andernfalls ist der gesuchte Zähler die nächstgrößere ganze Zahl, also jener Quotient um eins vermehrt, nachdem die folgenden Dezimalziffern weggelassen wurden. Der Grund dieser Regel ergibt sich so leicht aus den Bemerkungen am Ende des vorigen Artikels, dass eine ausführlichere Erklärung unnötig ist.

\textit{Beispiel.} Weiß man, dass die ersten beiden Ziffern der Mantisse eines Bruchs mit dem Nenner \(23\) gleich \(69\) sind, so haben wir das Produkt \(23.69=1587\). Entfernt man dessen letzte beiden Ziffern und addiert eins, so ergibt sich der gesuchte Zähler \(=16\).

\Needspace{7\baselineskip}
\Article{314}
Wir beginnen mit Brüchen, deren Nenner Primzahlen oder Primzahlpotenzen sind, und zeigen später, wie sich die übrigen darauf zurückführen lassen. Zunächst bemerken wir sofort, dass die Mantisse eines Bruchs \(\frac a{p^\mu}\), von dessen Zähler \(a\) wir stets voraussetzen, dass er durch die Primzahl \(p\) nicht teilbar ist, endlich ist und aus \(\mu\) Ziffern besteht, wenn \(p=2\) oder \(=5\). Im ersten Fall ist diese als ganze Zahl betrachtete Mantisse \(=5^\mu a\), im zweiten \(=2^\mu a\). Dies ist so offensichtlich, dass es keiner Erklärung bedarf.

Ist dagegen \(p\) eine andere Primzahl, so ist \(10^r a\) niemals durch \(p^\mu\) teilbar, wie groß auch \(r\) gewählt wird. Daher muss die Mantisse von \(F=\frac a{p^\mu}\) unendlich fortlaufen. Es sei \(10^e\) die kleinste Potenz von \(10\), die nach dem Modul \(p^\mu\) der Einheit kongruent ist. Man vergleiche Abschnitt III, wo wir gezeigt haben, dass \(e\) entweder der Zahl \((p-1)p^{\mu-1}\) gleich oder ein echter Teiler davon ist. Dann sieht man leicht, dass auch \(10^e a\) die erste Zahl in der Reihe \(10a,100a\), \(1000a\) usw. ist, die nach demselben Modul der Zahl \(a\) kongruent ist. Nun entstehen nach Art.\ \(312\) die Mantissen der Brüche \(\frac{10a}{p^\mu},\frac{100a}{p^\mu}\ldots\frac{10^e a}{p^\mu}\), indem man von der Mantisse
\SourceLeaf{389}
des Bruchs \(F\) die erste Ziffer, die ersten beiden \(\ldots e\) Ziffern der Reihe nach entfernt. Daher wiederholen sich in dieser Mantisse offenbar nach den ersten \(e\) Ziffern, und nicht früher, dieselben Ziffern. Diese ersten \(e\) Ziffern, deren unendliche Wiederholung die Mantisse bildet, können wir die \textit{Periode} dieser Mantisse oder des Bruchs \(F\) nennen. Offenbar ist die Länge der Periode, \textit{das heißt}, die Anzahl ihrer Ziffern (\(=e\)), vom Zähler \(a\) völlig unabhängig und nur durch den Nenner bestimmt. So ist \textit{zum Beispiel} die Periode von \(\frac1{11}\) gleich \(09\) und die von \(\frac37\) gleich \(428571\)\HistoricalNote{*)}{Der angesehene Robertson bezeichnet Anfang und Ende der Periode durch zwei über die erste und die letzte Ziffer gesetzte Punkte (\textit{Theory of circulating fractions, Philos. Trans.} \(1769\) \textit{S.} \(207\)), was wir hier nicht für nötig halten.}.

\Needspace{7\baselineskip}
\Article{315}
Kennt man also die Periode eines Bruchs, so kann seine Mantisse auf beliebig viele Ziffern fortgesetzt werden. Gilt ferner \(b\equiv10^\lambda a\; (\mathrm{mod}.\ p^\mu)\), so erhält man offenbar die Periode von \(\frac b{p^\mu}\), indem man die ersten \(\lambda\) Ziffern der Periode von \(F\), wobei wir \(\lambda<e\) voraussetzen dürfen, hinter die übrigen \(e-\lambda\) Ziffern stellt. Mit der Periode von \(F\) hat man somit zugleich die Perioden aller Brüche, deren Zähler den Zahlen \(10a,100a,1000a\) usw. nach dem Nenner \(p^\mu\) kongruent sind. So ergibt sich \textit{zum Beispiel}, da \(6\equiv3.10^2\; (\mathrm{mod}.\ 7)\), die Periode von \(\frac67\) sofort aus der von \(\frac37\) als \(857142\).

Ist nach dem Modul \(p^\mu\) die Zahl \(10\) eine Primitivwurzel, nach Art.\ \(57,89\), so lässt sich aus der Periode von \(\frac1{p^\mu}\) unmittelbar die Periode jedes anderen Bruchs \(\frac m{p^\mu}\) gewinnen, dessen Zähler \(m\) durch \(p\) nicht teilbar ist. Man entfernt links so viele Ziffern und fügt sie rechts wieder an, wie der Index von \(m\) Einheiten hat, wenn \(10\) als Basis gewählt wird. Daraus wird deutlich, warum in diesem Fall in Tafel I stets die Zahl \(10\) als Basis gewählt wurde; siehe Art.\ \(72\).

\begin{samepage}
Ist dagegen \(10\) keine Primitivwurzel, so lassen sich aus der Periode von \(\frac1{p^\mu}\) nur die Perioden derjenigen Brüche gewinnen, deren Zähler einer Potenz von \(10\) nach \(p^\mu\) kongruent sind. Es sei \(10^e\) die kleinste Potenz von \(10\), die nach \(p^\mu\) der Einheit kongruent ist, ferner \((p-1)p^{\mu-1}=ef\). Als Basis wähle man eine solche Primitivwurzel \(r\), dass der Index von \(10\) gleich \(f\) wird, nach Art.\ \(71\). In diesem System haben die Zähler der Brüche, deren Perioden aus der Periode von \(\frac1{p^\mu}\) gewonnen werden können, die Indizes \(f,2f,3f\ldots ef-f\). Ebenso gewinnt man aus der Periode von \(\frac r{p^\mu}\) die Perioden der Brüche, deren Zähler \(10r,100r,1000r\) usw. den Indizes \(f+1,2f+1,3f+1\) usw. entsprechen. Aus der Periode des Bruchs mit dem Zähler \(rr\), dessen Index \(2\) ist, gewinnt man die Perioden der Brüche, deren Zähler
\SourceLeaf{390}
die Indizes \(f+2,2f+2,3f+2\) usw. haben. Allgemein lassen sich aus der Periode des Bruchs mit dem Zähler \(r^i\) die Perioden der Brüche gewinnen, deren Zähler die Indizes \(f+i,2f+i,3f+i\) usw. haben. Daraus folgt leicht: Kennt man nur die Perioden der Brüche mit den Zählern \(1,r,rr,r^3\ldots r^{f-1}\), so lassen sich alle übrigen allein durch Umstellung nach folgender Regel ableiten. Der Index des Zählers \(m\) des gegebenen Bruchs \(\frac m{p^\mu}\) im System mit der Basis \(r\) sei \(=i\), wobei wir ihn kleiner als \((p-1)p^{\mu-1}\) voraussetzen. Durch Division durch \(f\) setze man \(i=\alpha f+\beta\), sodass \(\alpha,\beta\) positive ganze Zahlen oder auch \(0\) sind und \(\beta<f\) gilt. Dann entsteht die Periode von \(\frac m{p^\mu}\) aus der Periode des Bruchs mit dem Zähler \(r^\beta\), also \(1\), wenn \(\beta=0\), indem man die ersten \(\alpha\) Ziffern dieser Periode hinter die übrigen stellt. Für \(\alpha=0\) behält man also diese Periode selbst. Dies erklärt hinreichend, warum wir beim Aufstellen der Tafel I der in Art.\ \(72\) erläuterten Regel gefolgt sind.
\par\end{samepage}
